\section{Historical training and contextual comparisons}
\subsection{Domain adaptation and re-anchoring}

The studies collected here provide supporting evidence under their original
protocols, separately from the sealed RF experiment. Table~\ref{tab:history-ladder}
compares the raw-base, supervised, and earlier optimized checkpoints, retaining
the full declared draw budget for each. The earlier
\claim{DesignCount}-design evaluation provides broader five-family coverage:
its penalized endpoint rises from \claim{OverallSftPenalized} to
\claim{OverallGrpoPenalized}~MHz, with \claim{OverallImproved} of
\claim{DesignCount} designs improving. Feed-forward families show the most
consistent gains, while IIR is mixed and median remains slow in absolute terms.
Because its uncertainty analysis is historical and post-hoc, this study
supports the breadth of the findings but is not pooled with the sealed result.

\begin{table}[b]
\centering
\caption{Earlier five-family capability ladder. Utility is failure-penalized
timing-derived frequency; correctness denotes the sampled oracle verdict.}
\label{tab:history-ladder}
\small
\begin{tabular}{lrr}
\toprule
Checkpoint & Utility (MHz) & Oracle pass (\%) \\
\midrule
Raw base & \claim{BasePenalized} & \claim{BaseCorrectness} \\
SFT & \claim{OverallSftPenalized} & \claim{OverallSftCorrectness} \\
Earlier policy (v8) & \claim{OverallGrpoPenalized} & \claim{OverallGrpoCorrectness} \\
\bottomrule
\end{tabular}
\end{table}

% AUTO-GENERATED by paper/make_revision_evidence.py.
\begin{table}[b]
\centering
\caption{Historical MLP re-anchoring evaluation on its own complete draw set. This is separate from the sealed RF study and the earlier v8 evaluation.}
\label{tab:history-vnine}
\small
\begin{tabular}{@{}lrr@{}}
\toprule
Checkpoint & Utility (MHz) & Oracle pass (\%) \\
\midrule
Matched SFT & \revisionclaim{HistoryVnineSftUtility} & \revisionclaim{HistoryVnineSftPass} \\
Re-anchored policy (v9) & \revisionclaim{HistoryVninePolicyUtility} & \revisionclaim{HistoryVninePolicyPass} \\
\bottomrule
\end{tabular}

\end{table}

The historical re-anchored MLP policy (v9) raises utility from
\revisionclaim{HistoryVnineSftUtility} to
\revisionclaim{HistoryVninePolicyUtility}~MHz on its own matched SFT sample,
with \revisionclaim{HistoryVnineImproved} of
\revisionclaim{HistoryVnineDesigns} designs improving. Its mean gain is
\revisionclaim{HistoryVnineGain}~MHz. Oracle acceptance also increases
(Table~\ref{tab:history-vnine}). The earlier v8 and v9 evaluations use different
training budgets and generated samples. Their difference cannot isolate the
causal effect of re-anchoring, and training-only seed logs do not supply
additional held-out physical replications.

\subsection{Second-backbone study}
Qwen-Coder repeats the complete domain-SFT and policy-optimization recipe on
the earlier \claim{QwenDesignCount}-design,
\claim{QwenFamilyCount}-family universe. This is a separate earlier pipeline
experiment, with its own SFT initialization and reward configuration. The main
paper retains its overall and regime-level results because they provide
positive evidence beyond a single backbone. They are not a second test of the
canonical-feature RF package, and are not pooled with the sealed study.

\subsection{API prompting and HLS context}
The archived API experiment uses the same accelerator tasks and oracle, with
\claim{FrontierSamplesPerDesign} samples per prompt arm. HLS uses engineer-written
C++ and either no pragmas or engineer-selected pragmas. The reported RTL
comparison selects the fastest measured optimized candidate. These budgets,
inputs, and selection rules differ from the main policy-average comparison.
Within the API experiment, a timing-focused prompt raises conditional \fmax{} from
\claim{FrontierPlainConditional} to
\claim{FrontierFastConditional}~MHz while correctness changes from
\claim{FrontierPlainCorrectness}\% to
\claim{FrontierFastCorrectness}\%; its penalized endpoint is
\claim{FrontierFastPenalized}~MHz. This contrast illustrates why conditional
speed and per-draw utility need to be reported separately. It describes one API
configuration under a smaller sampling budget and does not establish a general
ranking of frontier models.

For selected implementation quality, the fastest measured optimized RTL
candidate reaches \claim{HlsExpertGeomean}$\times$ the geometric-mean
throughput of engineer-pragma HLS over \claim{HlsDesignCount} comparable
designs.  It reaches \claim{HlsNaiveGeomean}$\times$ no-pragma HLS.  Because HLS
starts from engineer-written C++ and the RTL value is selected rather than
policy-average, this comparison is context, not the causal baseline.
\subsection{Earlier general-RTL evaluation}
The earlier base-model comparison showed pass-at-one decreasing from
\claim{VerilogBasePassone}\% to \claim{VerilogSftPassone}\% after domain SFT.
That experiment uses different generation settings; its absolute values are
not comparable with the matched RF extension and are not pooled.  Together,
the two experiments identify specialization costs at different training stages,
without establishing performance on unseen industrial RTL.
